\section{Plano horizontal: junta 1}
\label{sec:ik_horizontal}

Pela cinemática direta (Equação~\ref{eq:xyz}, com $l = L_5$), o ponto $(x, y)$ é a soma do desvio lateral $L_5$ com a distância $d$ em frente, rodadas de $q_1$ em torno do eixo $z$ (Figura~\ref{fig:topo}):
\begin{equation*}
    x = L_5\cos q_1 - d\sin q_1, \qquad
    y = L_5\sin q_1 + d\cos q_1
\end{equation*}

\subsection{Distância em frente}

Uma rotação conserva as distâncias: como $L_5$ e $d$ são perpendiculares, a distância $r$ do alvo ao eixo de J1 satisfaz o teorema de Pitágoras, $x^2 + y^2 = L_5^2 + d^2$. Logo:
\begin{equation}
    r = \sqrt{x^2 + y^2}, \qquad d = \sqrt{r^2 - L_5^2}
    \label{eq:ik_d}
\end{equation}

Só existe solução se $r > L_5$: um alvo mais perto do eixo de J1 do que o desvio lateral está dentro do próprio braço. Toma-se a raiz positiva de $d$, que corresponde ao braço virado para o ponto, em frente da base.

\subsection{Ângulo da junta 1}

O alvo faz o ângulo $\operatorname{atan2}(y, x)$ com o eixo $x$. A rotação $q_1$ soma-se ao ângulo que o alvo faz com a direção do desvio, $\operatorname{atan2}(d, L_5)$, pelo que $\operatorname{atan2}(y, x) = q_1 + \operatorname{atan2}(d, L_5)$ e
\begin{equation}
    q_1 = \operatorname{atan2}(y, x) - \operatorname{atan2}(d, L_5)
    \label{eq:ik_q1}
\end{equation}

É por causa do desvio lateral $L_5$ que o segundo termo aparece. Sem ele, o plano do braço passaria pelo eixo de J1 e bastaria a fórmula da aula, $q_1 = \operatorname{atan2}(y, x)$, descontando apenas os $90^\circ$ que vêm de, com $q_1 = 0$, o braço apontar para $+y$ e não para $+x$.

\clearpage